\IfFileExists{stacks-project.cls}{%
\documentclass{stacks-project}
}{%
\documentclass{amsart}
}

% For dealing with references we use the comment environment
\usepackage{verbatim}
\newenvironment{reference}{\comment}{\endcomment}
%\newenvironment{reference}{}{}
\newenvironment{slogan}{\comment}{\endcomment}
\newenvironment{history}{\comment}{\endcomment}

% For commutative diagrams we use Xy-pic
\usepackage[all]{xy}

% We use 2cell for 2-commutative diagrams.
\xyoption{2cell}
\UseAllTwocells

% We use multicol for the list of chapters between chapters
\usepackage{multicol}

% This is generally recommended for better output
\usepackage{lmodern}
\usepackage[T1]{fontenc}

% For cross-file-references
\usepackage{xr-hyper}

% Package for hypertext links:
\usepackage{hyperref}

% For any local file, say "hello.tex" you want to link to please
% use \externaldocument[hello-]{hello}
\externaldocument[introduction-]{introduction}
\externaldocument[conventions-]{conventions}
\externaldocument[sets-]{sets}
\externaldocument[categories-]{categories}
\externaldocument[topology-]{topology}
\externaldocument[sheaves-]{sheaves}
\externaldocument[sites-]{sites}
\externaldocument[stacks-]{stacks}
\externaldocument[fields-]{fields}
\externaldocument[algebra-]{algebra}
\externaldocument[brauer-]{brauer}
\externaldocument[homology-]{homology}
\externaldocument[derived-]{derived}
\externaldocument[simplicial-]{simplicial}
\externaldocument[more-algebra-]{more-algebra}
\externaldocument[smoothing-]{smoothing}
\externaldocument[modules-]{modules}
\externaldocument[sites-modules-]{sites-modules}
\externaldocument[injectives-]{injectives}
\externaldocument[cohomology-]{cohomology}
\externaldocument[sites-cohomology-]{sites-cohomology}
\externaldocument[dga-]{dga}
\externaldocument[dpa-]{dpa}
\externaldocument[sdga-]{sdga}
\externaldocument[hypercovering-]{hypercovering}
\externaldocument[schemes-]{schemes}
\externaldocument[constructions-]{constructions}
\externaldocument[properties-]{properties}
\externaldocument[morphisms-]{morphisms}
\externaldocument[coherent-]{coherent}
\externaldocument[divisors-]{divisors}
\externaldocument[limits-]{limits}
\externaldocument[varieties-]{varieties}
\externaldocument[topologies-]{topologies}
\externaldocument[descent-]{descent}
\externaldocument[perfect-]{perfect}
\externaldocument[more-morphisms-]{more-morphisms}
\externaldocument[flat-]{flat}
\externaldocument[groupoids-]{groupoids}
\externaldocument[more-groupoids-]{more-groupoids}
\externaldocument[etale-]{etale}
\externaldocument[chow-]{chow}
\externaldocument[intersection-]{intersection}
\externaldocument[pic-]{pic}
\externaldocument[weil-]{weil}
\externaldocument[adequate-]{adequate}
\externaldocument[dualizing-]{dualizing}
\externaldocument[duality-]{duality}
\externaldocument[discriminant-]{discriminant}
\externaldocument[derham-]{derham}
\externaldocument[local-cohomology-]{local-cohomology}
\externaldocument[algebraization-]{algebraization}
\externaldocument[curves-]{curves}
\externaldocument[resolve-]{resolve}
\externaldocument[models-]{models}
\externaldocument[functors-]{functors}
\externaldocument[equiv-]{equiv}
\externaldocument[pione-]{pione}
\externaldocument[etale-cohomology-]{etale-cohomology}
\externaldocument[proetale-]{proetale}
\externaldocument[relative-cycles-]{relative-cycles}
\externaldocument[more-etale-]{more-etale}
\externaldocument[trace-]{trace}
\externaldocument[crystalline-]{crystalline}
\externaldocument[spaces-]{spaces}
\externaldocument[spaces-properties-]{spaces-properties}
\externaldocument[spaces-morphisms-]{spaces-morphisms}
\externaldocument[decent-spaces-]{decent-spaces}
\externaldocument[spaces-cohomology-]{spaces-cohomology}
\externaldocument[spaces-limits-]{spaces-limits}
\externaldocument[spaces-divisors-]{spaces-divisors}
\externaldocument[spaces-over-fields-]{spaces-over-fields}
\externaldocument[spaces-topologies-]{spaces-topologies}
\externaldocument[spaces-descent-]{spaces-descent}
\externaldocument[spaces-perfect-]{spaces-perfect}
\externaldocument[spaces-more-morphisms-]{spaces-more-morphisms}
\externaldocument[spaces-flat-]{spaces-flat}
\externaldocument[spaces-groupoids-]{spaces-groupoids}
\externaldocument[spaces-more-groupoids-]{spaces-more-groupoids}
\externaldocument[bootstrap-]{bootstrap}
\externaldocument[spaces-pushouts-]{spaces-pushouts}
\externaldocument[spaces-chow-]{spaces-chow}
\externaldocument[groupoids-quotients-]{groupoids-quotients}
\externaldocument[spaces-more-cohomology-]{spaces-more-cohomology}
\externaldocument[spaces-simplicial-]{spaces-simplicial}
\externaldocument[spaces-duality-]{spaces-duality}
\externaldocument[formal-spaces-]{formal-spaces}
\externaldocument[restricted-]{restricted}
\externaldocument[spaces-resolve-]{spaces-resolve}
\externaldocument[formal-defos-]{formal-defos}
\externaldocument[defos-]{defos}
\externaldocument[cotangent-]{cotangent}
\externaldocument[examples-defos-]{examples-defos}
\externaldocument[algebraic-]{algebraic}
\externaldocument[examples-stacks-]{examples-stacks}
\externaldocument[stacks-sheaves-]{stacks-sheaves}
\externaldocument[criteria-]{criteria}
\externaldocument[artin-]{artin}
\externaldocument[quot-]{quot}
\externaldocument[stacks-properties-]{stacks-properties}
\externaldocument[stacks-morphisms-]{stacks-morphisms}
\externaldocument[stacks-limits-]{stacks-limits}
\externaldocument[stacks-cohomology-]{stacks-cohomology}
\externaldocument[stacks-perfect-]{stacks-perfect}
\externaldocument[stacks-introduction-]{stacks-introduction}
\externaldocument[stacks-more-morphisms-]{stacks-more-morphisms}
\externaldocument[stacks-geometry-]{stacks-geometry}
\externaldocument[moduli-]{moduli}
\externaldocument[moduli-curves-]{moduli-curves}
\externaldocument[examples-]{examples}
\externaldocument[exercises-]{exercises}
\externaldocument[guide-]{guide}
\externaldocument[desirables-]{desirables}
\externaldocument[coding-]{coding}
\externaldocument[obsolete-]{obsolete}
\externaldocument[fdl-]{fdl}
\externaldocument[index-]{index}

% Theorem environments.
%
\theoremstyle{plain}
\newtheorem{theorem}[subsection]{Theorem}
\newtheorem{proposition}[subsection]{Proposition}
\newtheorem{lemma}[subsection]{Lemma}

\theoremstyle{definition}
\newtheorem{definition}[subsection]{Definition}
\newtheorem{example}[subsection]{Example}
\newtheorem{exercise}[subsection]{Exercise}
\newtheorem{situation}[subsection]{Situation}

\theoremstyle{remark}
\newtheorem{remark}[subsection]{Remark}
\newtheorem{remarks}[subsection]{Remarks}

\numberwithin{equation}{subsection}

% Macros
%
\def\lim{\mathop{\mathrm{lim}}\nolimits}
\def\colim{\mathop{\mathrm{colim}}\nolimits}
\def\Spec{\mathop{\mathrm{Spec}}}
\def\Hom{\mathop{\mathrm{Hom}}\nolimits}
\def\Ext{\mathop{\mathrm{Ext}}\nolimits}
\def\SheafHom{\mathop{\mathcal{H}\!\mathit{om}}\nolimits}
\def\SheafExt{\mathop{\mathcal{E}\!\mathit{xt}}\nolimits}
\def\Sch{\mathit{Sch}}
\def\Mor{\mathop{\mathrm{Mor}}\nolimits}
\def\Ob{\mathop{\mathrm{Ob}}\nolimits}
\def\Sh{\mathop{\mathit{Sh}}\nolimits}
\def\NL{\mathop{N\!L}\nolimits}
\def\CH{\mathop{\mathrm{CH}}\nolimits}
\def\proetale{{pro\text{-}\acute{e}tale}}
\def\etale{{\acute{e}tale}}
\def\QCoh{\mathit{QCoh}}
\def\Ker{\mathop{\mathrm{Ker}}}
\def\Im{\mathop{\mathrm{Im}}}
\def\Coker{\mathop{\mathrm{Coker}}}
\def\Coim{\mathop{\mathrm{Coim}}}

% Boxtimes
%
\DeclareMathSymbol{\boxtimes}{\mathbin}{AMSa}{"02}

%
% Macros for moduli stacks/spaces
%
\def\QCohstack{\mathcal{QC}\!\mathit{oh}}
\def\Cohstack{\mathcal{C}\!\mathit{oh}}
\def\Spacesstack{\mathcal{S}\!\mathit{paces}}
\def\Quotfunctor{\mathrm{Quot}}
\def\Hilbfunctor{\mathrm{Hilb}}
\def\Curvesstack{\mathcal{C}\!\mathit{urves}}
\def\Polarizedstack{\mathcal{P}\!\mathit{olarized}}
\def\Complexesstack{\mathcal{C}\!\mathit{omplexes}}
% \Pic is the operator that assigns to X its picard group, usage \Pic(X)
% \Picardstack_{X/B} denotes the Picard stack of X over B
% \Picardfunctor_{X/B} denotes the Picard functor of X over B
\def\Pic{\mathop{\mathrm{Pic}}\nolimits}
\def\Picardstack{\mathcal{P}\!\mathit{ic}}
\def\Picardfunctor{\mathrm{Pic}}
\def\Deformationcategory{\mathcal{D}\!\mathit{ef}}
\usepackage{cleveref}
\crefname{originallemma}{original lemma}{original lemmas}
\Crefname{originallemma}{Original lemma}{Original lemmas}
\makeatletter
\newlabel{lemma-split-into-connected}{{09WS}{source}{}{sourcepointer-1}{}}
\newlabel{lemma-split-into-connected@cref}{{[originallemma][0][]09WS}{[1][0][]source}{}{}{}}
\makeatother
\crefname{originalsection}{original section}{original sections}
\Crefname{originalsection}{Original section}{Original sections}
\makeatletter
\newlabel{section-limit-sets}{{002U}{source}{}{sourcepointer-2}{}}
\newlabel{section-limit-sets@cref}{{[originalsection][0][]002U}{[1][0][]source}{}{}{}}
\makeatother
\title{The colimit condition preserves injections of set diagrams but not of abelian group diagrams}
\author{The Stacks Project Authors}
\date{stacks-project:categories.tex:local-snapshot-20261012; GNU FDL 1.2 or later}
\begin{document}
\maketitle
\noindent\textit{Editorial note.} Counted span is the author's own proof environment, source lines 2586-2666, following the statement at 2565-2584.\par\medskip
\noindent\textit{Editorial note.} Part (1) covers diagrams of sets; part (2) is the counterexample for abelian groups, both inside the same span.\par\medskip
\noindent\textit{Allowed standard prerequisites (editorial).}\begin{itemize}
\item Decomposition into connected components. Lemma lemma-split-into-connected of this chapter, invoked at proof lines 2588-2593. Writes I as a coproduct of nonempty connected index categories satisfying the same condition, reducing (1) to the connected case since colimits over a coproduct decompose.
\item Colimits of set-valued diagrams as a quotient description. Section section-limit-sets of this chapter, used without further mention at 2600-2601. The zig-zag chain description of equality in a colimit of sets; the induction runs on chain length.
\end{itemize}\medskip
\ifdefined\hypertarget\hypertarget{sourcepointer-1}{}\fi
\noindent\textit{Source reference.} Original lemma 09WS: sources/stacks-categories/pinned-original-source.tex, source lines 2538--2538. Navigation pointer only: the counted proof cites this statement by \textbackslash{}ref as a standing input; the deduction the counted proof has to make is carried out inside the counted span, and no part of that deduction is delegated to this pointer.\par\medskip
\ifdefined\hypertarget\hypertarget{sourcepointer-2}{}\fi
\noindent\textit{Source reference.} Original section 002U: sources/stacks-categories/pinned-original-source.tex, source lines 1807--1807. Navigation pointer only: the counted proof cites this statement by \textbackslash{}ref as a standing input; the deduction the counted proof has to make is carried out inside the counted span, and no part of that deduction is delegated to this pointer.\par\medskip
\begin{lemma}
\label{lemma-preserve-injective-maps}
Let $\mathcal{I}$ be an index category, i.e., a category. Assume
that for every solid diagram
$$
\xymatrix{
x \ar[d] \ar[r] & y \ar@{..>}[d] \\
z \ar@{..>}[r] & w
}
$$
in $\mathcal{I}$ there exist an object $w$ and dotted arrows
making the diagram commute. Then
\begin{enumerate}
\item an injective morphism $M \to N$ of diagrams of sets over
$\mathcal{I}$ gives rise to an injective map $\colim M_i \to \colim N_i$
of sets,
\item in general the same is not the case for diagrams of abelian
groups and their colimits.
\end{enumerate}
\end{lemma}

\begin{proof}
If $\mathcal{I}$ is the empty category, then the lemma is true.
Thus we may assume $\mathcal{I}$ is nonempty. In this case
we can write $\mathcal{I} = \coprod \mathcal{I}_j$ where each
$\mathcal{I}_j$ is nonempty and satisfies the same property, see
Lemma \ref{lemma-split-into-connected}. Since
$\colim_\mathcal{I} M = \coprod_j \colim_{\mathcal{I}_j} M|_{\mathcal{I}_j}$
this reduces the proof of (1) to the connected case.

\medskip\noindent
Assume $\mathcal{I}$ is connected and $M \to N$ is injective, i.e.,
all the maps $M_i \to N_i$ are injective.
We identify $M_i$ with the image of $M_i \to N_i$, i.e.,
we will think of $M_i$ as a subset of $N_i$.
We will use the description of the colimits given in
Section \ref{section-limit-sets} without further mention.
Let $s, s' \in \colim M_i$ map to the same element of $\colim N_i$.
Say $s$ comes from an element $m$ of $M_i$ and $s'$ comes from an
element $m'$ of $M_{i'}$. Then we can find a sequence
$i = i_0, i_1, \ldots, i_n = i'$ of objects of $\mathcal{I}$
and morphisms
$$
\xymatrix{
&
i_1 \ar[ld] \ar[rd] & &
i_3 \ar[ld] & &
i_{2n-1} \ar[rd] & \\
i = i_0 & &
i_2 & &
\ldots & &
i_{2n} = i'
}
$$
and elements $n_{i_j} \in N_{i_j}$ mapping to each other under
the maps $N_{i_{2k-1}} \to N_{i_{2k-2}}$ and $N_{i_{2k-1}}
\to N_{i_{2k}}$ induced from the maps in $\mathcal{I}$ above
with $n_{i_0} = m$ and $n_{i_{2n}} = m'$. We will prove by induction
on $n$ that this implies $s = s'$. The base case $n = 0$ is trivial.
Assume $n \geq 1$. Using the assumption on $\mathcal{I}$
we find a commutative diagram
$$
\xymatrix{
& i_1 \ar[ld] \ar[rd] \\
i_0 \ar[rd] & & i_2 \ar[ld] \\
& w
}
$$
We conclude that $m$ and $n_{i_2}$ map to the same element of $N_w$
because both are the image of the element $n_{i_1}$.
In particular, this element is an element $m'' \in M_w$ which
gives rise to the same element as $s$ in $\colim M_i$.
Then we find the chain
$$
\xymatrix{
&
i_3 \ar[ld] \ar[rd] & &
i_5 \ar[ld] & &
i_{2n-1} \ar[rd] & \\
w & &
i_4 & &
\ldots & &
i_{2n} = i'
}
$$
and the elements $n_{i_j}$ for $j \geq 3$ which has a smaller length
than the chain we started with. This proves the induction step and the
proof of (1) is complete.

\medskip\noindent
Let $G$ be a group and let $\mathcal{I}$ be the one-object category with
endomorphism monoid $G$. Then $\mathcal{I}$ satisfies the condition stated
in the lemma because given $g_1, g_2 \in G$ we can find $h_1, h_2 \in G$
with $h_1 g_1 = h_2 g_2$. A diagram $M$ over $\mathcal{I}$ in
$\textit{Ab}$ is the same thing as an abelian group $M$ with $G$-action
and $\colim_\mathcal{I} M$ is the coinvariants $M_G$ of $M$.
Take $G$ the group of order $2$ acting trivially on $M = \mathbf{Z}/2\mathbf{Z}$
mapping into the first summand of
$N = \mathbf{Z}/2\mathbf{Z} \times \mathbf{Z}/2\mathbf{Z}$
where the nontrivial element of $G$ acts by
$(x, y) \mapsto (x + y, y)$. Then $M_G \to N_G$ is zero.
\end{proof}

\end{document}
